% Partitiones Oratoriae (partitiones-oratoriae) --- English translation
% Translated by Claude (Anthropic) from the Latin (Perseus).
% Style guide: STYLE.md.

\ciceroSection{1} \textsc{Marcus.} I should like, my father, to hear from you in Latin those things on the theory of speaking which you taught me in Greek—if only you have the leisure, and if you are willing. \textsc{Cicero.} Is there anything, my Marcus, that I should prefer to your being as learned as possible? As for leisure, in the first place it is at its fullest, since at last the chance to leave Rome has been given me; and in the second, I would gladly set those pursuits of yours ahead of even my weightiest occupations.

\ciceroSection{2} \textsc{Marcus.} Do you wish, then, that just as you are accustomed to question me in Greek in due order, so I should in turn question you in Latin on the same matters? \textsc{Cicero.} By all means, if you like. For in this way I shall both understand that you remember what you have learned, and you will hear in due order the things you seek.

\ciceroSection{3} \textsc{Marcus.} Into how many parts is the whole science of speaking to be divided? \textsc{Cicero.} Three. \textsc{Marcus.} Tell me which? \textsc{Cicero.} First, into the faculty of the orator himself; next, into the speech; then, into the question. \textsc{Marcus.} In what does the faculty itself lie? \textsc{Cicero.} In matter and in words. But both matter and words must be discovered and arranged. Properly, however, discovery is said of the matter, expression of the words. As for arrangement, although it is common to both, it is nonetheless referred to discovery. The voice, the movement, the countenance, and all delivery are the companions of expression; and the guardian of all these things is memory.

\ciceroSection{4} \textsc{Marcus.} And the speech—how many parts has it? \textsc{Cicero.} Four. Two of them serve to instruct in the matter, the narrative and the proof; two to drive the feelings, the introduction and the peroration. \textsc{Marcus.} And the question—what parts has it? \textsc{Cicero.} The indefinite, which I call a general inquiry, and the definite, which I name a case.

\ciceroSection{5} \textsc{Marcus.} Since, then, discovery is the orator's first task, what will he seek? \textsc{Cicero.} How to find the way to produce conviction in those whom he wishes to persuade, and the way to bring emotion into their minds. \textsc{Marcus.} By what means is conviction produced? \textsc{Cicero.} By arguments, which are drawn from topics either inherent in the matter itself or brought in from outside. \textsc{Marcus.} What do you call topics? \textsc{Cicero.} Those in which arguments lie hidden. \textsc{Marcus.} What is an argument?

\ciceroSection{6} \textsc{Cicero.} A plausible thing discovered to produce conviction. \textsc{Marcus.} How, then, do you divide those two kinds? \textsc{Cicero.} Those that are thought to be without art I call external, such as testimonies; the inherent ones cleave to the matter itself. \textsc{Marcus.} What kinds of testimony are there? \textsc{Cicero.} The divine and the human. The divine is such as oracles and auspices, such as prophecies and the responses of priests, soothsayers, and diviners; the human is what is regarded from authority, from intention, from speech whether free or extracted under pressure—and in it are contained things written, agreed, promised, sworn, and obtained by examination.

\ciceroSection{7} \textsc{Marcus.} What are the ones you call inherent? \textsc{Cicero.} The relations fixed in the matters themselves: such as definition; such as the contrary; such as the things that belong either to the thing itself, or to its contrary, or to things like or unlike, in agreement or in disagreement; such as the things that are as it were conjoined, or as it were at war with one another; such as the causes of the matters under discussion; such as the outcomes of those causes—that is, the things produced from the causes; such as distributions, the genus of the kinds or the kinds of the genus; such as the first beginnings of things and, as it were, their forerunners, in which there is something of argument; such as the comparisons of things—what is greater, what equal, what less—in which either the natures of things or their capacities are weighed against each other.

\ciceroSection{8} \textsc{Marcus.} Shall we, then, take arguments from all of those topics? \textsc{Cicero.} On the contrary, we shall search and seek from all of them, but we shall apply judgment, so as always to throw out the trivial ones, and sometimes also to pass over the commonplace and the unnecessary. \textsc{Marcus.} Since you have answered about conviction, I want to hear about emotion. \textsc{Cicero.} You ask in the right place, but what you wish will be more fully unfolded when I come to the theory of the speech itself and of the questions.

\ciceroSection{9} \textsc{Marcus.} What follows, then? \textsc{Cicero.} When you have discovered, you must arrange; and in the indefinite question the order is roughly the same that I set out for the topics, while in the definite there must be added also those things that pertain to the emotions. \textsc{Marcus.} How, then, do you unfold these matters? \textsc{Cicero.} I have common precepts for producing conviction and for stirring feeling. Since conviction is a firm belief, while emotion is a rousing of the mind toward pleasure or toward distress, toward desire or toward fear—for so many are the kinds of emotion, with several parts to each single kind—I fit the whole arrangement to the end of the question. For in the indefinite question the end is conviction, in the case both conviction and emotion. And so, when I have spoken of the case, in which the indefinite question is contained, I shall have spoken of both. \textsc{Marcus.} What, then, have you to say about the case?

\ciceroSection{10} \textsc{Cicero.} That it is distinguished by the kind of its audience. For the one who hears is either a mere listener, or an arbiter—that is, one who decides on the matter and the verdict; so that the one who hears either is meant to be delighted, or to settle something. And he settles either about things past, as a juror does, or about things to come, as a senator does; hence these three kinds—the judicial, the deliberative, and the display kind—the last of which, because it is applied chiefly to acts of praise, has by now taken its proper name from that: panegyric.

\ciceroSection{11} \textsc{Marcus.} What aims does the orator set before himself in those three kinds? \textsc{Cicero.} Delight in panegyric; in a trial, either the harshness or the clemency of the juror; in persuasion, either the hope or the dread of the man deliberating. \textsc{Marcus.} Why, then, do you set out the kinds of cases at this point? \textsc{Cicero.} So that I may fit the theory of arrangement to the end of each.

\ciceroSection{12} \textsc{Marcus.} In what way, pray? \textsc{Cicero.} Because in those speeches whose end is delight, the orders of arrangement are various. For we either keep the gradations of time or the distributions of kinds; or we mount up from the lesser to the greater, or glide down from the greater to the lesser; or we mark these off with an uneven variety, when we interweave the small with the great, the simple with the conjoined, the obscure with the lucid, the cheerful with the sad, the incredible with the plausible—all of which belong to panegyric.

\ciceroSection{13} \textsc{Marcus.} And in deliberation, what do you look to? \textsc{Cicero.} Introductions either short, or often none at all; for those who deliberate are prepared, for their own sakes, to listen. Nor, indeed, must much often be narrated; for narrative is of things past or present, while persuasion is of things to come. For this reason the whole speech must be directed to conviction and to emotion.

\ciceroSection{14} \textsc{Marcus.} And in trials, what is the arrangement? \textsc{Cicero.} Not the same for the accuser and for the defendant. For the accuser pursues the order of events and arranges the several arguments, as it were spears, in his hand; he sets them forth forcefully, drives his conclusion home sharply, confirms it with records, decrees, and testimonies, and dwells rather carefully on the single points, using those precepts of speech that serve to rouse the feelings; in the rest of his speech he digresses a little from the course of his pleading, and grows more vehement in the peroration. For his aim is to make the juror angry.

\ciceroSection{15} \textsc{Marcus.} What must be done on the defendant's side? \textsc{Cicero.} Everything quite the reverse. The introduction must be taken up to win goodwill; the narratives must be either cut away, the ones that wound, or left out entirely, if they are wholly damaging; the supports laid down by the other side to produce conviction must be either dissolved one by one, or obscured, or buried under digressions; while the perorations must be directed to pity. \textsc{Marcus.} Can we, then, always keep the order of arrangement we wish? \textsc{Cicero.} By no means; for the ears of the audience govern the prudent and far-sighted orator: what they reject must be changed.

\ciceroSection{16} \textsc{Marcus.} Set out next what the precepts are of the speech itself and of the words. \textsc{Cicero.} There is one kind of expression that flows of its own accord; another that is turned and altered. The first power lies in simple words, the second in conjoined ones. Simple words must be discovered, the conjoining of them must be arranged. And of simple words, some are native, some invented. The native are those given by the senses; the invented are those made from these and coined anew, whether by likeness, or by imitation, or by inflection, or by the joining together of words.

\ciceroSection{17} There is, besides, this distinction in words: one of nature, the other of treatment. Of nature, in that some are more sonorous, grander, lighter, and in a certain way more polished, others the contrary; of treatment, when one takes either the proper terms of things, or terms added to a name, or new ones, or archaic ones, or ones modified and bent in a certain way by the orator—such as those that are transferred or altered, or those that we as it were misapply, or those that we obscure, or that we raise to an incredible pitch, or that we adorn more wondrously than the habit of common speech allows.

\ciceroSection{18} \textsc{Marcus.} I have it about simple words; now I ask about their conjoining. \textsc{Cicero.} Certain rhythms must be observed in the conjoining, and a sequence of words. The rhythms the ears themselves measure, so that you neither fall short of completing in words what you have set out to say, nor run over. The sequence, in turn, so that the speech is not thrown into confusion in its genders, numbers, tenses, persons, and cases. For just as in simple words what is not Latin is to be censured, so in conjoined words what is not consequent.

\ciceroSection{19} Common to simple and conjoined words are these five, as it were, lights: the lucid, the brief, the plausible, the brilliant, the agreeable. Lucidity will be achieved by familiar, proper words set in order, whether rounded off in a period, or by a pause, or by the clipping of words; obscurity, on the other hand, comes either by length or by contraction of the speech, by ambiguity, or by the inflection and alteration of words. Brevity is brought about by simple words, by saying each thing once and once only, by serving no end but that of speaking lucidly. The plausible kind of speech is one that is not too trim and polished, that has authority and weight in its words, with thoughts either weighty or fitted to the beliefs and habits of men.

\ciceroSection{20} A speech is brilliant if the words are set down chosen for their weight, and transferred, and heightened, and joined to a name, and doubled, and ones of like meaning, and not at odds with the very action and imitation of things. For this is the part of a speech that sets the matter almost before the eyes; for it is chiefly that sense which is touched, yet the others too, and most of all the mind itself, can be moved. But the things said about lucid speech all fall under this brilliant kind as well. For the brilliant is something considerably more than the lucid: by the one we are brought to understand, by the other we seem to see.

\ciceroSection{21} The agreeable kind of speaking will arise, first, from elegance and from the charm of words that are sonorous and light; then from a conjoining that has neither harsh collisions nor disjointed and gaping gaps, that is rounded off not by a long winding but by a compass fitted to the breath of the voice, and that has a likeness and evenness of words—when, from words drawn out of contraries, word answers to word, like to like, and when words referred to the same term, and twinned and doubled, or even repeated still oftener, are set down, and the construction of words is now coupled by conjunctions, now as it were loosened by dissolutions.

\ciceroSection{22} A speech will also become agreeable when you say something unusual, or unheard of, or new. Whatever is wonderful gives delight too; and most of all that speech moves the hearer which rouses some emotion in the mind, and which marks out the lovable character of the orator himself—a character expressed either by the stamp of his own judgment and a humane and generous spirit, or by an inflection of speech, when, for the sake either of magnifying another or of diminishing himself, some things seem to be said by the orator and others to be meant, and this seems done out of courtesy rather than out of pretense. But there are many precepts of agreeableness that make a speech either more obscure or less plausible. And so here too we must ourselves judge what the case demands.

\ciceroSection{23} \textsc{Marcus.} It remains, then, that you speak of the turned and altered speech. \textsc{Cicero.} That whole kind lies in the changing of words, which is handled in simple words in such a way that the speech is either expanded out of a single word or contracted into a single word: out of a word, when a proper word, or one of like meaning, or a coined word, is drawn out into several words; out of a phrase, when either a definition is recalled to a single word, or words brought in are removed, or roundabout phrases are straightened, or by conjunction one word is made out of two.

\ciceroSection{24} In conjoined words, however, a threefold change can be applied—not of the words, but only of the order: so that, when something has been said once straight out, as nature herself has carried it, the order is reversed, and the same thing is said as it were upside down and backward, and then the same again broken up and intermingled. And the exercise of expression turns most of all upon this whole kind of inversion.

\ciceroSection{25} \textsc{Marcus.} Delivery follows next, I suppose. \textsc{Cicero.} It does; and delivery indeed is the thing the orator must vary most, in keeping with the weight of both the matter and the words. For it makes a speech lucid, brilliant, plausible, and agreeable, not by words, but by variety of voice, by movement of body, by countenance—all of which will avail most if they accord with the kind of speech and follow its force and its variety.

\ciceroSection{26} \textsc{Marcus.} Does anything at all remain to you concerning the orator himself? \textsc{Cicero.} Nothing, surely, except memory, which is in a certain way the twin of writing, and, in a dissimilar kind, very like it. For just as writing consists of marks for letters and of that on which the marks themselves are imprinted, so the working of memory uses topics as if they were wax, and arranges images on them as if they were letters.

\ciceroSection{27} \textsc{Marcus.} Since, then, the whole faculty of the orator has been set out, what have you to say about the precepts of the speech? \textsc{Cicero.} That it has four parts, of which the first and the last serve to stir the feeling of the mind—for the mind must be roused in the openings and the perorations—while the second, the narrative, and the third, the proof, produce conviction for the speech. But amplification, although it has its own proper place, often indeed the first, and almost always the last, must nonetheless be applied throughout the rest of the course of the speech, and most of all when something has been either confirmed or refuted. And so it avails very greatly for conviction too; for amplification is a kind of forceful argumentation: as that serves the purpose of instruction, so this serves the purpose of stirring.

\ciceroSection{28} \textsc{Marcus.} Go on, then, in order, and unfold for me those four parts. \textsc{Cicero.} I will, and I shall begin first with the introductions, which are drawn either from persons or from the matters themselves. They are taken up for the sake of three things: that we may be heard with goodwill, with understanding, and with attention. The first ground for these lies in our own persons, and in those of the arbiters and the adversaries; from these the beginnings of winning goodwill are gathered, whether from our own deserts, or from our standing, or from some kind of virtue, and most of all from generosity, sense of duty, justice, and good faith; while contrary things are turned against the adversaries, and we point to some bond with those who decide, whether of relationship or of hope; and if any hatred or offense has been brought upon us, it must be removed or lessened, whether by dissolving it, or by extenuating it, or by counterbalancing it, or by pleading against it.

\ciceroSection{29} That we may be heard intelligently, and likewise attentively, we must begin from the matter itself. The hearer learns most easily, and grasps what is at issue, if at the outset you take in the kind and nature of the case, if you define, if you divide, and if you neither hinder his understanding by a confusion of the parts nor his memory by their number. And what will presently be said about a lucid narrative may rightly be transferred to this place as well.

\ciceroSection{30} That we may be heard attentively, we shall achieve by one of three means: for we shall set out something great, or something necessary, or something bound up with the very people before whom the matter is being tried. Let this too stand among the precepts: that if ever the occasion itself, or the matter, or the place, or someone's intervention, or an interruption, or some saying of the adversary's—and especially in the peroration—has given us a chance to say something fitting to the moment, we should not let it slip. And of what we shall say in its own place about amplification, much can be carried over to the precepts on introductions.

\ciceroSection{31} \textsc{Marcus.} What, then, must be observed in the narrative? \textsc{Cicero.} Since the narrative is the unfolding of the matter, and as it were the seat and foundation on which conviction is established, those things must be most of all observed in it which are observed in nearly all the other parts of speaking as well; and these are partly necessary, partly assumed for the sake of adornment. For that we narrate lucidly and plausibly is necessary, but we add charm besides.

\ciceroSection{32} So for narrating lucidly we shall call again upon those same earlier precepts of making plain and shedding light, among which is that brevity which is so very often praised in the narrative, of which I spoke above. It will be plausible if the things narrated agree with the persons, the times, the places; if a cause is set down for each deed and outcome; if they seem to be told on testimony, if joined with the opinion and authority of men, with the law, with custom, with religion; if the honesty of the narrator is signified, his age, his reputation, the truthfulness of his speech and the trustworthiness of his life. The narrative is charming when it has wonders, suspense, unexpected outcomes, stirrings of feeling set between, exchanges between persons, griefs, angers, fears, joys, desires. But let us now press on to what remains.

\ciceroSection{33} \textsc{Marcus.} Surely what follows are the things that pertain to producing conviction. \textsc{Cicero.} It is so: and these are divided into proof and refutation. For in proving we wish to make good our own contentions, in refuting to confute the contrary ones. Since, then, everything that comes into dispute is questioned either as to whether it is or is not, or as to what it is, or as to what its quality is, in the first conjecture prevails, in the second definition, in the third reasoning. \textsc{Marcus.} I have that distribution. Now I ask for the topics of conjecture.

\ciceroSection{34} \textsc{Cicero.} It rests wholly upon what is probable and upon the proper marks of things. But let us, for teaching's sake, call probable that which usually happens so—as that youth is rather prone to lust—and a proper mark that argument which is never otherwise and declares a thing for certain, as smoke declares fire. Probabilities will be found in the parts and as it were the limbs of the narrative. They lie in persons, in places, in times, in deeds, in outcomes, in the natures of the matters and affairs themselves.

\ciceroSection{35} In persons there are considered first the marks of nature—of health, figure, strength, age, of men and women; and these indeed in the body; but as to the mind, either how it is disposed, by virtues and vices, by skills and ineptitudes, or how it is stirred, by desire, fear, pleasure, distress. And these are the things considered in nature; in fortune there are birth, friendship, children, kinsfolk, relations by marriage, wealth, honours, powers, riches, freedom, and the things contrary to these.

\ciceroSection{36} In places there are both those natural features—whether they are by the sea or remote from it, level or mountainous, smooth or rough, healthy or pestilent, shaded or sunlit—and those that are matters of chance: whether tilled or untilled, frequented or deserted, built up or empty, obscure or made famous by the traces of deeds done, consecrated or profane.

\ciceroSection{37} In times the present, the past, and the future are discerned; and within these very ones, the long-past, the recent, the now-pressing, what is to be a little later or at some time. There lie in times also those features that mark as it were the nature of the season—as winter, as summer, or the times of the year, as the month, as the day, as night, hour, weather, which are natural; and the matters of chance, such as sacrifices, festival days, weddings.

\ciceroSection{38} Now deeds and outcomes are matters either of design or of imprudence, which lies either in accident or in some disturbance of mind: in accident, when a thing fell out otherwise than was thought; in disturbance, when either forgetfulness, or error, or some cause of fear or desire moved a man. Under imprudence necessity too is to be set down. Of good things and evil there are three kinds; for they can lie either in the mind, or in the body, or outside. With this matter, then, laid down for the argument, all the parts must be surveyed in the mind, and from each single one a conjecture must be drawn toward the point at issue.

\ciceroSection{39} There is also another kind of arguments which is taken from the traces of the deed, such as a weapon, blood, a cry heard, faltering, a change of colour, inconsistent speech, trembling of the accused or of others, anything that can be grasped by the senses; and likewise, if something was prepared beforehand, if it was shared with another, if afterward it was seen, heard, disclosed.

\ciceroSection{40} Probabilities, again, partly carry weight singly by their own force, and partly, even though they seem slight in themselves, accomplish much when they are heaped together. And among these probabilities there lie sometimes also the sure and proper marks of things. But what produces the greatest conviction of likeness to truth is, first, an example, then a likeness of the matter brought in; sometimes too a tale, even if it is past belief, nevertheless stirs men.

\ciceroSection{41} \textsc{Marcus.} What, then, is the method and the way of definition? \textsc{Cicero.} There is no doubt of this, that a definition is set forth by the genus and by a certain property, or even by a frequency of common features, out of which what is proper may shine clear. But since over the proper features a great disagreement usually arises, we must often define from contraries, often too from things unlike, often from things equal. For this reason descriptions also are often apt in this kind, and an enumeration of consequences; and above all the unfolding of the term and the name carries weight.

\ciceroSection{42} \textsc{Marcus.} Already set forth, I think, are nearly all the things that are questioned about the deed and about the naming of the deed. So then there remain, surely, those things which—when both the deed is established and its name—are called into doubt as to what their quality is? \textsc{Cicero.} It is just as you say. \textsc{Marcus.} What, then, are the parts in that kind? \textsc{Cicero.} That the deed was done lawfully, or for the sake of warding off or avenging a wrong, or in the name of duty, or chastity, or religion, or country, or finally out of necessity, ignorance, accident.

\ciceroSection{43} For the deeds that are done by an impulse and disturbance of mind, without reason, have no defences against the charge in lawful trials, though they may have them in free debates. In this kind, in which the quality is questioned, it is usual to ask, out of the dispute, whether the thing was done lawfully or not; and the argument of this must be drawn from the description of the topics.

\ciceroSection{44} \textsc{Marcus.} Come then, since you had divided the conviction of a speech into proof and refutation, and the one has been spoken of, set forth now about refuting. \textsc{Cicero.} Either the whole of what the adversary has assumed in his argument must be denied, if you can show it to be feigned or false; or the things assumed as probabilities must be confuted: first, that doubtful things have been assumed for certain; next, that the same could be said even of things plainly false; then, that from the things he has assumed there does not follow what he wishes. And it must come to this: that as the single points, so the whole, will be broken. Examples too must be recalled in which, in a like dispute, no credit was given; and the condition of a common peril must be lamented, if the life of the innocent is to be laid open to the wits of slanderous men.

\ciceroSection{45} \textsc{Marcus.} Since I now have where the things that pertain to conviction are found, I await how the several points are to be handled in the speaking. \textsc{Cicero.} You seem to be asking for argumentation, which is the unfolding of an argument; and this, taken from the topics that have been set forth, must be worked out and distinguished lucidly. \textsc{Marcus.} That very thing is plainly what I am after.

\ciceroSection{46} \textsc{Cicero.} It is, then, as was said above: argumentation is the unfolding of an argument; but it is worked out when you have assumed things either not doubtful or probable, from which you bring about that which seems doubtful or less probable in itself. Of arguing there are two kinds, of which the one looks directly toward conviction, the other bends itself toward the stirring of feeling. Directly, then, when one has set forth something to be proved and has assumed the things on which it might rest, and, these being confirmed, has referred himself back to what was proposed and concluded. But that other argumentation as it were backwards and contrariwise first assumes the things it wishes and confirms them, then casts at the end, after the feelings are stirred, that which was to have been proposed.

\ciceroSection{47} There is also that variety in arguing, and that not unpleasant differentiation, when we put a question to ourselves, or inquire, or command, or wish—which, along with several others, are ornaments of the thought. We shall be able to avoid sameness by not always beginning from the proposition; and if we do not confirm everything by disputing, but now and then set down briefly the things that are clear enough; and what is brought about from these, if it is plain, we shall not always need to conclude.

\ciceroSection{48} \textsc{Marcus.} What of those things that are called artless—which you said a while ago were assumed—do they in any way, in any place, stand in need of art? \textsc{Cicero.} They do indeed stand in need of it; nor are they called artless because they are so, but because the orator's art does not give birth to them—yet, brought to it from outside, it handles them all the same with art, and most of all in the matter of witnesses.

\ciceroSection{49} For both about the whole class of witnesses one must often say how weak a thing it is, and that arguments are proper to the matters, but testimonies to men's wills; and one must use examples, where no credit has been given to witnesses; and about individual witnesses, whether they are by nature empty, frivolous, marked with disgrace, or driven by hope, by fear, by anger, by pity, or led by reward, by favour; and they are to be compared with the higher authority of witnesses to whom, even so, no credit was given.

\ciceroSection{50} Often, too, examinations under torture must be resisted, because many, fleeing pain, have very often lied amid the torments and chosen rather to die by confessing a falsehood than to suffer by denying the truth; many also have set their own life at nought, that they might free those who were dearer to them than they were to themselves; while others, either by the constitution of their body or by a habituation to pain or by the fear of execution and death, have borne up under the force of the torments; and others have lied against those they hated. And these things must be made good by examples.

\ciceroSection{51} Nor is it obscure that, since on either side there are examples and likewise topics for forming a conjecture, in contrary cases contrary things must be assumed. And there comes into play, besides, a certain other method in the matter of witnesses and of examinations under torture. For often the things that have been said, if they were said ambiguously or inconsistently or past belief, or even said differently by different men, are subtly refuted.

\ciceroSection{52} \textsc{Marcus.} There remains for you the last part of the speech, which is set in the peroration, about which I should very much like to hear. \textsc{Cicero.} The unfolding of the peroration is easier. For it is divided into two parts, amplification and recapitulation. The proper place for heightening is here, in the peroration; and in the very course of the speech turns toward amplification are given, when some matter has been confirmed or refuted.

\ciceroSection{53} Amplification, then, is a kind of weightier affirmation, which by the stirring of feeling wins conviction in the speaking. It is worked out both by the kind of words and by the matter. Words are to be set down which have the force of shedding light and are not at odds with usage—weighty, full, sonorous, compounded, coined, surnamed, not common, exalted, and above all transferred. These in single words; but in continuous passages the loosened kind, which are spoken without conjunction, that they may seem more.

\ciceroSection{54} Words also heighten when brought back, repeated, doubled, and those that rise step by step from lowlier words to higher; and altogether the speech that is, as it were, natural and not made plain, but stuffed with weighty words, is the better fitted for heightening. These things, then, in the words, to which there is suited an action of the voice and a gesture matching and apt to stir feeling. But both in the words and in the action the case must be weighed, and one must act according to the matter. For because these things seem utterly absurd when they are weightier than the case bears, it must be judged with care what befits each occasion.

\ciceroSection{55} The amplification of the matter, in turn, is taken from all those same topics from which were taken the things that were said toward conviction: and most of all there prevail definitions massed together, and a heaping-up of consequences, and the clashing of contrary, unlike, and mutually warring things, and the causes and the things that have arisen from causes, and most of all likenesses and examples; feigned persons too, and let mute things at last speak; and altogether those things are to be brought in, if the case allows, which are held great—of which there are two kinds.

\ciceroSection{56} For some things seem great by nature, others by use: by nature, such as the heavenly things, the divine, those whose causes are obscure, the marvels there are on earth and in the world—from which, and from things like them, if you attend, very many means of heightening are at hand; by use, the things that seem to men more strongly to profit or to harm, of which there are three kinds for amplifying. For men are moved either by attachment, as to the gods, to country, to parents; or by love, as for brothers, for spouses, for children, for intimates; or by honour, as for the virtues, and most of all those that bear upon the fellowship of men and upon generosity. From these are taken both exhortations to hold fast to such things, and hatreds are kindled against those by whom they have been violated, and pity is born.

\ciceroSection{57} The proper place for amplification lies in such goods as these, when they are lost or in danger of being lost. For nothing is so pitiable as one made wretched out of a happy state. And the whole of what stirs feeling is this: from what fortune a man falls, and from the love of what people he is torn away, what he loses or has lost, what evils he is in or is about to be in — all this is expressed briefly. For a tear dries quickly, especially over another man's misfortunes. Nor should anything be too minutely worked out in amplification, since all painstaking detail is a thing of small things; but this place demands the grand.

\ciceroSection{58} It must now be left to judgement what kind of amplification we are to use in each case. For in those cases which are dressed up for the sake of delight, the topics to be handled are those that can rouse expectation, wonder, and pleasure; in exhortations, the enumerations of goods and evils, and examples, count for most. In trials, what bears on anger is generally for the accuser, what bears on pity for the defendant; though sometimes the accuser too ought to stir pity, and the defender, anger.

\ciceroSection{59} The enumeration remains: for the panegyrist never necessary, for the adviser not often, for the accuser more often than for the defendant. Of its uses there are two: if you should distrust the memory of those before whom you plead, whether through the lapse of time or the length of the speech, or when, the proofs of the speech being massed together and set out briefly, the case will have greater force.

\ciceroSection{60} The defendant must use it more rarely, because contraries have to be set out, the dissolving of which will not shine in brevity, while their stings will prick. But in the enumeration this is to be avoided: that the display of memory should not seem to be undertaken in a childish way. He will escape this who does not run back over every smallest point, but, touching on the single matters briefly, grasps the very weight of the things.

\ciceroSection{61} \textsc{Marcus.} Since you have spoken both of the orator himself and of the speech, set out for me now that last division of the question which you proposed as the third. \textsc{Cicero.} There are, as I said at the start, two kinds of question, of which the one, bounded by times and persons, I call a case; the other, unbounded, marked off by no persons and no times, I call a thesis. But the thesis is, as it were, a part of the controverted case; for the unbounded is contained within the bounded, and yet all things are referred back to it.

\ciceroSection{62} For this reason let us speak first of the thesis: of which there are two kinds, the one of knowledge — its end is understanding, as whether the senses are true; the other of action, which is referred to the doing of something, as if it were asked by what duties friendship is to be cultivated. Again, of the former there are three kinds: whether a thing is or is not, and what it is, and of what quality it is. Whether it is or is not: as whether right lies in nature or in custom; what it is, thus: whether that is right which is useful to the greater part; of what quality it is, thus: whether to live justly is useful.

\ciceroSection{63} Of action, however, there are two kinds: one for the pursuing of something or the avoiding of it, as by what means you may attain glory, or in what way envy may be shunned; the other, which is referred to some advantage and use, as in what way the commonwealth is to be administered, or in what way one must live in poverty.

\ciceroSection{64} And again, out of the inquiry of knowledge — where it is asked whether a thing is or is not, or has been, or is to be — one kind of question is whether something can be brought about, as when it is asked whether anyone at all can be perfectly wise; another, in what way each thing comes to be, as by what means virtue is begotten, whether by nature or by reason or by practice. To this kind belong all those questions in which, as in obscure and natural inquiries, the causes and reasons of things are unfolded.

\ciceroSection{65} But of that kind in which it is asked what the thing in question is, there are two species: in the one it must be argued whether a thing is different or the same, as obstinacy and perseverance; in the other a description of some class and, as it were, an image of it must be drawn out, as what the miser is like, or who the proud man is.

\ciceroSection{66} In the third kind, in which it is asked of what quality a thing is, we must speak either of the honourable or of the advantageous or of the equitable. Of the honourable, thus: whether it is honourable to undergo danger or unpopularity for a friend. Of the advantageous, thus: whether it is advantageous to engage in administering the commonwealth. Of the equitable, thus: whether it is fair to set friends before kinsmen. And within this same kind, in which it is asked of what quality a thing is, there arises another kind of argument. For it is not simply asked what is honourable, what advantageous, what fair, but also, by comparison, what is more honourable, what more advantageous, what fairer, and even what is most honourable, most advantageous, most fair; of which kind are these questions, what dignity of life is the most excellent. And all those I have mentioned belong to knowledge.

\ciceroSection{67} Those of action remain: of which the one is the kind of giving precepts, which bears on the rule of duty, as in what way parents are to be honoured; the other is for the calming of minds and their healing by speech, as in consoling griefs, in repressing anger, in soothing fear, or in lessening desire. To this kind the contrary is the kind of argument directed at those same movements of feeling, which has often to be done in amplifying a speech — whether to engender or to rouse them.

\ciceroSection{68} And this is more or less the division of theses. \textsc{Marcus.} I have grasped it; but I ask what the method is of discovering and arranging matter in these. \textsc{Cicero.} What? Do you suppose it is any other than the same one already set out — that all things are drawn for conviction and for discovery from the same topics? And the method of arrangement, set out under other heads, will be carried over here unchanged. \textsc{Marcus.} The whole division of theses being known, then, the kinds and precepts of cases remain for us.

\ciceroSection{69} \textsc{Cicero.} Just so. And of these the form is twofold: the one pursues the delight of the ears, the other aims to win its right, to prove, to bring about what it is doing — from which the whole struggle is undertaken. And so the former is called panegyric: which, though it can be a broad and very varied kind, we single out from it the one part that we undertake for the praising of famous men and the censure of the wicked. For there is no kind of speaking that can be either richer in matter for speaking, or more useful to states, or in which the orator is more engaged in the knowledge of the virtues and vices. The remaining kind of cases turns either on provision for the time to come or on the disputing of the time past; of which the one belongs to deliberation, the other to trials.

\ciceroSection{70} Out of this division three kinds of cases have emerged: one, which has been named from the better part of praise, panegyric; the second, deliberation; the third, that of trials. For which reason let us discuss the first first, if it please you. \textsc{Marcus.} It does indeed please me. \textsc{Cicero.} Then I shall set out briefly the methods of praising and blaming — which avail not only for speaking well but also for living honourably — and I shall begin from the openings both of praise and of blame.

\ciceroSection{71} For surely all things bound up with virtue are to be praised, and all things bound up with vices are to be blamed. For which reason the end of the one is honour, of the other disgrace. This kind of discourse is achieved by narrating and setting forth deeds — which, without any argumentations, is suited rather to handling the movements of feeling gently than to producing or confirming conviction. For doubtful things are not made firm, but those things which are certain, or are laid down as certain, are amplified. For which reason from what has been said before the precepts both of narrating and of amplifying will be drawn.

\ciceroSection{72} And since in these cases the whole method is referred almost wholly to the pleasure of the hearer and to delight, we shall have to use in them ornamented speech, and the distinctions of single words, which have the most charm — this comes about if we frequently use words coined or old or transferred — and the very construction of the words, so that like answers to like and similar to similar, with contraries, with words doubled, with periods rhythmically rounded off, not to the likeness of verse, but to fill out the sense of the ears by a certain fitting measure, as it were, of the words.

\ciceroSection{73} And those ornaments of matter too must be brought in more frequently: whether such things as are wonderful, or unforeseen, or signified by portents, prodigies, and oracles, or such as will seem to have befallen the man we are speaking of by divine and fated agency. For every expectation of the hearer, and every wonder, and every unexpected outcome, hold some pleasure in the hearing.

\ciceroSection{74} But since goods and evils lie in three kinds — those outside us, those of the body, and those of the mind — let the external ones come first, which are drawn from family; and this being praised briefly and modestly, or, if it be infamous, passed over, if humble, either passed by or touched on so as to increase the glory of the man you praise, next, if the matter allows, we must speak of his fortunes and resources. Afterwards of the goods of the body; among which, that which most signifies virtue, as it were, beauty is the most easily praised.

\ciceroSection{75} Then we must come to deeds, whose arrangement is threefold: for either the order of times must be kept, or whatever is most recent must be told first, or many and various deeds must be sorted into the proper kinds of virtues. But this place of the virtues and vices, ranging most widely, will now, out of many and various discussions, be confined within a certain narrow and brief compass.

\ciceroSection{76} The force of virtue, then, is twofold: for virtue is discerned either in knowledge or in action. For that which is called prudence, that which is called shrewdness, and that which by the weightiest name is called wisdom, is strong in knowledge alone. But that which is praised for the governing of desires and the ruling of the movements of the mind has its function in acting; and its name is temperance. And that prudence, in one's own affairs, is wont to be called domestic; in public affairs, civic.

\ciceroSection{77} Temperance, likewise, is distributed into one's own affairs and into common ones, and is discerned in two ways in matters of advantage: both by not coveting things that are absent, and by abstaining from things that are within one's power. In matters of disadvantage it is likewise twofold: for that which stands against coming evils is named fortitude; that which bears and endures what is already present, patience. But that which embraces these under one kind is called greatness of mind; to which belongs liberality in the use of money, and at the same time loftiness of mind in taking on disadvantages and especially injuries — and everything of that kind that is grave, settled, not turbulent.

\ciceroSection{78} But the part that is set in community is called justice — and this, towards the gods, religion; towards parents, piety; in the common phrase, goodness; in matters entrusted, good faith; in the restraint of the mind, gentleness; in goodwill it is named friendship. And these virtues, indeed, are discerned in acting. But there are others, as it were the handmaids and companions of wisdom: of which the one distinguishes and judges what is true and false in argument, and what follows from given premises — a virtue which lies wholly in the reason and science of disputing; the other is oratorical.

\ciceroSection{79} For eloquence is nothing else than wisdom speaking copiously; which, drawn from the same source as that wisdom employed in argument, is richer and broader, and better suited to the movements of feeling and to the senses of the multitude. But the keeper of all the virtues, fleeing disgrace and most of all attaining praise, is modesty. And these, indeed, are more or less certain dispositions of the mind, so affected and so constituted that each is marked off from the others by its own proper kind of virtue; and from these, in whatever way a thing is done, so must it of necessity be honourable and supremely worthy of praise.

\ciceroSection{80} But there are also certain other formed dispositions of the mind, fashioned and prepared, as it were, towards virtue by right pursuits and arts: as, in one's private affairs, the pursuits of letters, of numbers and sounds, of measurement, of the stars, of horses, of hunting, of arms; in common affairs, inclinations towards the cultivation of some one kind of virtue in particular, or towards devotion to divine matters, or to the loving of parents, friends, and guests with particular and signal devotion. And these, indeed, belong to the virtues. But of the vices the kinds are contrary.

\ciceroSection{81} They must be discerned with care, lest those vices deceive us which seem to imitate virtue. For malice imitates prudence; brutality, temperance, in the spurning of pleasures; pride, greatness of mind, in puffing oneself up too far, and contemptuousness in despising honours; profligacy imitates liberality; rashness, fortitude; brutal hardness, patience; harshness, justice; superstition, religion; softness of mind, gentleness; timidity, modesty; the wrangling and catching at words imitates that prudence in argument, and a certain empty fluency of speech imitates this oratorical force. And to good pursuits there seem similar those that are excessive in the same kind.

\ciceroSection{82} For which reason the whole force of praising and blaming will be drawn from these parts of the virtues and vices; but in the whole, as it were, fabric of the speech these things will above all have to be made plain: in what way each man was born, in what way reared, in what way trained and shaped in his ways; and if anything great or incredible has befallen anyone, and especially if it can seem to have befallen by divine agency; then what each has felt, said, and done will be fitted to the kinds of virtue that have been set out, and from those topics of discovery the causes of things, and outcomes, and consequences will be sought. Nor indeed ought the death of those whose life is praised to be passed over in silence, if only there be something to remark either in the very manner of the death or in those things which have followed after death.

\ciceroSection{83} \textsc{Marcus.} I have received and learned these things in brief — not only in what way I might praise another, but also in what way I might strive to be justly praised myself. Let us see, then, next, what path and what precepts we are to hold in speaking an opinion. \textsc{Cicero.} In deliberation, then, the end is advantage, to which all things in giving counsel and speaking an opinion are so referred that these are the first things the adviser or dissuader must look to: what can or cannot be done, and what is or is not necessary. For if something cannot be brought about, deliberation is removed, however advantageous it may be; and if something is necessary — and that is necessary without which we cannot be safe or free — it is to be set before the remaining honourable considerations in civil reasoning, and before advantages.

\ciceroSection{84} But when it is asked what can be done, it must also be seen how easily it can be done. For things that are extremely difficult are often to be held just as if they could not be brought about. And when we attend to necessity, even if something will not seem necessary, it will yet have to be seen how great it is. For what is of very great moment is often held in place of the necessary.

\ciceroSection{85} And so, since this kind of case consists of persuasion and dissuasion, the man who urges a course is set a single line of argument: if it is both advantageous and possible, let it be done. The man who dissuades has a double one: the first, if it is not advantageous, that it not be done; the second, if it cannot be done, that it not be undertaken. Thus the man who urges must establish both points; for the man who dissuades it is enough to demolish either one.

\ciceroSection{86} Since, then, all deliberation turns on these two heads, let us speak first of the advantageous, which lies in the distinguishing of goods and evils. Of goods, some are necessary, such as life, chastity, freedom; such as children, spouses, brothers, parents. Others are not necessary; and of these, some are to be sought for their own sake, such as those that lie in our duties and virtues, while others are to be sought because they bring about some advantage, such as wealth and abundance.

\ciceroSection{87} Of those that are sought for their own sake, some are sought for their honourableness itself, some for some advantage. For their honourableness, those that proceed from the virtues of which I spoke a little earlier, which are praiseworthy of themselves; for some advantage, those things to be sought which lie in the goods of the body or of fortune. Of these, some are as it were joined with a certain honourableness, such as honour and glory; others are separate, such as strength, beauty, and health; such as noble birth, riches, and a body of clients.

\ciceroSection{88} There is also a certain matter, so to speak, that underlies honourableness, which is seen most of all in friendships. Friendships are discerned by affection and by love. For the reverence we owe the gods, and likewise our parents and our country, and those men who excel either in wisdom or in resources, is usually referred to affection; whereas spouses and children and brothers and the others whom intimacy and familiarity have joined to us, though they too are held by affection itself, are nonetheless held above all by love. Since, then, these things are goods, it is easy to understand what their contraries are.

\ciceroSection{89} Now if we could always hold to what is best, we should certainly have no great need of deliberation, since at least these matters are plain. But because, owing to circumstances, which have the greatest force, it very often happens that the advantageous contends with the honourable, and the clash of these things for the most part produces our deliberations, lest either the opportune be abandoned for the sake of dignity or the honourable for the sake of advantage, let us apply our precepts to the unravelling of this difficulty.

\ciceroSection{90} And since the speech is to be adapted not only to the truth but also to the opinions of those who listen, let us first understand this: that there are two kinds of men, the one untaught and rude, which always sets the advantageous before the honourable; the other cultivated and refined, which prefers dignity to everything. And so to this latter kind praise, honour, glory, good faith, justice, and every virtue are held out, but to the former gain, profit, and returns. Indeed pleasure too, which is most hostile to virtue and, by counterfeiting the nature of the good, corrupts it deceitfully, which the more brutish a man is the more keenly he pursues, and sets before not only honourable but even necessary things, is often quite worthy to be praised in persuasion, when you are giving counsel to that kind of men.

\ciceroSection{91} And this too must be observed: how much more men flee evils than pursue goods. For they do not seek the honourable so much as they shun the base. For who ever sought honour, glory, praise, or any distinction so eagerly as he flees disgrace, infamy, contumely, and dishonour? The keen pain of these things is grave witness that the human race, born for honour, has been corrupted by bad rearing and crooked opinions. And so in exhorting and persuading our aim will be this: to teach by what means we can attain goods and avoid evils;

\ciceroSection{92} but among well-bred men we shall speak chiefly of praise and of honour, and shall treat above all those kinds of virtue which are engaged in guarding and increasing the common good of mankind. But if we speak among the untaught and unskilled, let returns, profits, pleasures, and the avoidance of pains be brought forward; let contumely and disgrace be added too. For no one is so rude that, if honourableness itself moves him less, contumely and dishonour do not nonetheless move him greatly. And so what bears on the advantageous will be found from what has been said;

\ciceroSection{93} but what can be brought about — in which it is also customary to ask how easily it can be done and how far it serves — must be seen above all from the causes that bring about each thing. Now there are several kinds of causes. For some accomplish a thing of themselves, others bring some force toward accomplishing it. And so let the former be called the effecting causes, and let the rest be set in that class such that without them the thing cannot be brought about.

\ciceroSection{94} But an effecting cause is in one case absolute and complete in itself, in another a thing that helps toward something and is, as it were, a partner in producing the effect; and the force of this kind is various and often either greater or less, so that even the one which has the greatest force is often called the sole cause. There are also other causes which are called effecting either on account of a beginning or on account of an outcome. But when it is asked what is best to do, either the advantageous or the hope of accomplishing it impels men's minds to assent.

\ciceroSection{95} And since we have now spoken of the advantageous, let us speak of the means of accomplishment. In this whole class we must ask with whom and against whom, at what time, in what place, and with what resources of arms, money, allies, or of those things that bear on accomplishing each end, we may be able to act. And it is not only the things that are at our disposal that must be seen, but also those that stand against us. And if from the comparison our side will be the more favourable, it must not only be made convincing that what we urge can be done, but care must also be taken that those things appear easy, favourable, and welcome. For those who dissuade, on the other hand, either the advantage must be shaken or the difficulties of accomplishment brought out — and not from other precepts, but from these same topics of persuasion.

\ciceroSection{96} Let each of the two, however, have for amplification a store of examples, whether recent ones, so that they may be the better known, or old ones, so that they may carry more authority; and let him be especially practised in this kind of thing, so as to be able to set either the advantageous and necessary above the honourable, or these above those. For stirring men's minds, if they are to be roused, those sentiments will profit most which bear on satisfying desires, or glutting hatred, or avenging wrongs. But if they are to be restrained, they must be admonished concerning the uncertain state of fortune and the doubtful outcomes of things to come, and concerning the keeping of their own fortunes, if these are favourable, but if adverse, concerning the danger. And these indeed are the topics of the peroration.

\ciceroSection{97} But the introductions in delivering an opinion ought to be brief. For the orator comes not as a suppliant before a judge, but as one who exhorts and urges to action. He ought therefore to set out in what spirit he speaks, what he wants, and on what matters he is about to speak, and to exhort his hearers, speaking briefly, to give him a hearing. The whole speech, moreover, ought to be plain and weighty, and more adorned with thoughts than with words.

\ciceroSection{98} \textsc{Marcus.} I have now learned the topics of panegyric and of persuasion. Now I await those that suit the courts; and I think this one kind remains for us. \textsc{Cicero.} You understand rightly. And the end of this kind is fairness; which is regarded not simply, but sometimes from comparison, as when there is a dispute over who is the truest accuser, or when possession of an inheritance is sought without a law or without a will. In such cases it is asked what is fairer and what is fairest; and for such cases the means of argument are sought from those topics of fairness of which we shall presently speak.

\ciceroSection{99} And even before the trial there is usually a contest over the very setting up of the trial, when it is asked whether the man who brings the action has a right of action, or whether he has it yet, or whether it has now ceased to be his, or whether the action lies under that law, in these words. And these matters, even if they have not been contested or decided or settled before the case comes into court, nonetheless often carry very great weight in the trials themselves, when it is said: You have claimed too much; you have claimed too late; the claim was not yours; not against me, not under this law, not in these words, not in this court.

\ciceroSection{100} The class of these cases is grounded in the civil law, which is grounded in the law or custom governing private and public affairs; and the knowledge of it, neglected by most orators, seems to us necessary for speaking. For this reason — concerning the setting up of actions, the taking up or undertaking of trials, the raising of objections to the unfairness of an action, the securing of fairness — because these matters are generally of such a kind that, although they often slip down into the trial itself, they nonetheless seem to call for treatment before the trial, I set them a little apart from the trials, more by the time of pleading than by any difference of kind. For all matters disputed concerning the civil law or concerning the equitable and good fall under that form of cases in which it is in question what sort of thing something is, of which we are about to speak, and which consists chiefly in fairness and in law.

\ciceroSection{101} In all cases, then, there are three stages, of which some one must be taken — if you cannot take more — for resistance. For you must take your stand either so as to deny that the thing in question was done, or, if you confess it was done, to deny that it has the force, and is the thing, with which your adversary charges you, or, if there can be no dispute either about the deed or about the name of the deed, to deny that what you are accused of is such as he says, and to defend that what you did was right and ought to be granted you.

\ciceroSection{102} Thus that first stand, the clash, as it were, with the adversary, is to be handled by a kind of conjecture; the second by the definition or description and shaping of a word; the third by argument from what is fair and right and humane, to win pardon. And since the defender must always not only resist at some fixed stand, whether by denial or by definition or by setting fairness against the charge, but must also supply the ground of his refusal, that first stand has its ground in the very denial and disavowal of the unjust charge, the negation of the deed; the second, in the fact that what the adversary lays down in the word is not in the thing; the third, in his defending as right what he confesses to have been done without any dispute over the name.

\ciceroSection{103} Then against each ground the accuser must set that which, were it not in the accusation, the case could not exist at all. And so let the things that are thus brought to bear be called the supports of the cases; though those that are brought against the grounds of the defence do not hold the cases together any more than the grounds of the defence themselves do. But for the sake of distinction let us call ground that which is brought by the defendant to refuse and to ward off the charge, without which he would have nothing to defend; and let us call buttress that which, on the other side, is brought to overthrow the ground, without which the accusation could not stand.

\ciceroSection{104} Now from the clash and, as it were, the collision of ground and buttress a certain question arises, which I call the point at issue; in which it is usually asked what comes to trial and what is disputed. For the first contention of the adversaries has a diffuse question — as in conjecture, Did Decius take the money; in definition, Did Norbanus diminish the majesty of the state; in fairness, Did Opimius rightly kill Gracchus. These questions, which have their first contention from accusing and resisting, are broad, as I said, and diffuse. The contention of grounds and buttresses brings the point at issue into a narrow compass. In conjecture there is none. For no one can, or ought, or is wont to render a ground for what he denies was done. And so in these cases the first question and the last point at issue are one and the same.

\ciceroSection{105} But in those cases where it is said: He did not diminish the majesty of the state by his too turbulent handling of the matter of Caepio; for it was the people's just resentment that then roused that violence, not the tribune's action; and the majesty of the state, since it is a certain greatness of the Roman people in the keeping of its power and right, was rather increased than diminished — and where it is met with this: Majesty lies in the dignity of the empire and of the name of the Roman people, and he diminishes it who by the violence of a mob calls the commonwealth to sedition — there will arise this point at issue: Did he diminish the majesty of the state who, with the will of the Roman people, carried through by violence a thing welcome and fair?

\ciceroSection{106} But in those cases where something is defended as rightly done or as a deed that ought to be conceded, when the ground of the deed has been supplied — as by Opimius: I acted by right, for the safety of all and the preservation of the commonwealth, and Decius replies: Not even the most criminal citizen could you by any right kill without trial — there arises this point at issue: Could he rightly, for the safety of the commonwealth, kill without condemnation a citizen who was overturning the state? Thus those points at issue which arise in these controversies, which are marked by fixed persons and times, become indefinite again when persons and times are stripped away, and are once more brought back to the form and method of general inquiries.

\ciceroSection{107} But among the weightiest buttresses these too must be placed: any contraries drawn from the written text of a law, or of a will, or of the very words of the court order, or of some stipulation or guarantee, that are set against the defence. Yet not even this kind comes up in those cases that are contained under conjecture. For what is denied to have been done cannot be confuted by a written text. Nor does it come into a definition by the nature of the writing itself. For even if some word of the writing must be defined as to what force it has — as when it is asked from a will what "the store" is, or from a law about an estate what "the fixtures and fittings" are — it is not the nature of the writing but the interpretation of the word that makes the controversy.

\ciceroSection{108} But when either more things than one are signified by the writing on account of the ambiguity of a word or of words, so that the opposing party may draw the meaning of the writing whither it serves and pleases him; or, if the writing is not ambiguous, so that one may either lead the will and intent away from the words of the writer, or defend oneself by another writing of contrary tenor on the same matter — then the point at issue arises out of the contention over the writing: so that in ambiguities it is disputed what is most signified; in the contention of the letter and the intent, which of the two the judge should rather follow; in conflicting writings, which is the more to be approved.

\ciceroSection{109} Now when the point at issue has been established, the orator ought to have his thesis, into which all the arguments drawn from the topics of discovery may be gathered. And although this is enough for one who sees what lies hidden in each topic, and who keeps those topics noted like so many treasuries of arguments, we shall nonetheless touch on the things that are proper to particular cases.

\ciceroSection{110} In conjecture, then, when the defendant is in denial, the accuser has these two things first — but I call him "accuser" for every claimant and plaintiff, for these same kinds of controversies can be raised in cases even without an accusation — but these two things are first for him: the cause and the outcome. By "cause" I mean the ground of accomplishing, by "outcome" that which has been accomplished. And the division of causes itself was set out a little before among the topics of persuasion.

\ciceroSection{111} For the things that were taught about future time in the taking of counsel — for what reason they would seem to have either advantage or the means of accomplishment — these same things he who argues about a deed will have to gather, so that he may show for what reason the things were both advantageous to the man he accuses and could have been accomplished by him. Conjecture is set in motion by advantage, if the man accused is said to have acted either in hope of goods or in fear of evils; and this is the sharper, the greater those things are set down to be in either kind. The emotions too bear on the cause of the deed:

\ciceroSection{112} if there was fresh anger, if old hatred, if a passion for revenge, if the pain of a wrong, if a craving for honour, for glory, for power, for money, if fear of danger, if debt, if straitened family means; if the man is bold, if reckless, if cruel, if without self-command, if heedless, if unwise, if in love, if of disturbed mind, if drunk; if he acted with hope of accomplishing it, or in the belief that he could conceal it, or, if it were brought to light, that he could ward off the charge, or break through the danger, or put it off to a distant time; or if the penalty of the court was lighter than the reward of the deed; or if the pleasure of the crime was greater than the pain of condemnation.

\ciceroSection{113} For it is roughly by these things that the suspicion of a deed is confirmed, when both motives of will and the means are found in the accused. In the will, the advantage is sought in the gaining of some benefit and the avoiding of some harm, so that either hope or fear may seem to have driven him, or some sudden stirring of the mind, which drives a man into wrongdoing even more swiftly than a calculation of advantage. So let this be what is said about the motive.

\ciceroSection{114} \textsc{Marcus.} I have it; and I ask what those consequences are which you said are produced by the motives. \textsc{Cicero.} Certain attendant signs of the past, the imprinted traces, as it were, of the deed; which indeed stir suspicion most of all and are, so to speak, the silent testimonies of crimes, and the more weighty for this reason, that whereas the motives seem to incriminate and accuse, in common, all who had any interest in the matter, these consequences touch precisely the very men who are accused: as a weapon, as a footprint, as bloodstains, as something seized which appears to have been carried off or snatched away, as an inconsistent answer, as faltering, as stumbling in speech, as being seen with someone from whom suspicion arises, as being seen in the very place where the crime was done, as pallor, as trembling, as something written or sealed or deposited. For such things as these are what make a charge suspicious, whether in the deed itself, or even before it was done, or afterward.

\ciceroSection{115} And if these are lacking, one will nonetheless have to rely upon the motives themselves and upon the means of carrying it out, with that common argument added: that the man was not so deranged as to be unable either to escape or to conceal the evidence of the deed, that he should have laid himself so open as to leave room for a charge. On the other side stands the common topic, that audacity is joined with rashness, not with prudence.

\ciceroSection{116} There follows then that topic for amplification: that one ought not to wait until he confesses; that offences are convicted by arguments; and here too examples will be brought forward.

\ciceroSection{117} And so much for arguments. But if there shall also be the resource of witnesses, the first kind to be praised is the very fact, and one must say that, lest the accused be caught by arguments, he by his own caution brought it about that he could not escape the witnesses; then let the witnesses be praised one by one—and what their praiseworthy qualities are has been said—; then that, even though an argument be strong, since it is nonetheless often false, one may rightly withhold belief; whereas a good and steadfast man, free of fault, cannot fail to be believed without fault in the judge. And further, if the witnesses shall be obscure or of slight standing, it will have to be said that conviction is not to be weighed by fortune, or that those are the most trustworthy witnesses to any matter who can most easily know the thing in question. But if examinations under torture that have been held, or the demand that they be held, will aid the case, the first thing to be confirmed will be the kind of examination; one must speak of the force of pain, of the opinion of our forefathers, who, had they not approved the whole practice, would certainly have rejected it;

\ciceroSection{118} of the institutions of the Athenians, of the Rhodians, the most learned of men, among whom even—what is most harsh—free men and citizens are tortured; of the institutions, too, of our own most prudent men, who, although they were unwilling to have slaves examined against their masters, nonetheless judged that an examination should be held concerning incest, and concerning the conspiracy that was carried out in my consulship. The argument by which men are accustomed to weaken examinations under torture is also to be derided, and to be called rehearsed and childish. Then conviction must carefully be produced that the examination was conducted diligently and without bias, and that what was said under it must be weighed by arguments and conjecture. And these, roughly, are the limbs of the prosecution.

\ciceroSection{119} The defence, on the other hand, has first the weakening of the motives: that there were none, or none so great, or none belonging to him alone, or that he could have obtained the same end more conveniently; or that he is not of such character, not of such a life; or that there were no stirrings of the mind, or none so ungovernable. He will employ the weakening of the means, moreover, if he shall show that strength, or spirit, or resources, or wealth were wanting; or an unfavourable time, or an unsuitable place, or many bystanders, none of whom he could trust; or that he was not so open as to undertake what he could not conceal, nor so deranged as to despise penalties and the courts.

\ciceroSection{120} The consequences, too, he will dispel by setting forth that those are not sure indications of the deed, since they could follow even when no crime has been committed; and he will take his stand on the particulars, and defend them either as proper to things he himself will say were done, rather than to the crime, or as shared by him with the accuser and as bound to count for his safety rather than against it; and he will refute the whole class of witnesses and examinations under torture, and, where he can, the particulars, drawing on the topics of rebuttal of which I spoke before.

\ciceroSection{121} The introductions to these cases will be set down by the accuser so as to rouse hostility through suspicion, and the common danger of treachery will be proclaimed, and minds will be stirred to give heed. By the defendant, on the other hand, there will be put forward a complaint of a charge trumped up and of suspicions raked together, and the accuser's snares, and likewise the common danger; and minds will be drawn toward pity, and the goodwill of the judges gathered in moderation. The accuser's narrative will be, as it were, a suspicious unfolding of the business limb by limb, with all the arguments scattered through it and the defences obscured. The defender, with the arguments of suspicion either passed over or obscured, will have to narrate the outcomes and accidents of the events themselves.

\ciceroSection{122} In confirming our own argumentations and weakening the contrary ones, the accuser will often have to rouse the emotions, the defendant to soften them. And this, indeed, each must do most of all in the peroration: the one by the massing of his arguments and a wholesale heaping-up; the other, if he shall have plainly unfolded his case by refutation, by an enumeration of how he has dispelled each point, and at the last by an appeal to pity.

\ciceroSection{123} \textsc{Marcus.} I think I now see how conjecture is to be handled. Now let us hear about definition. \textsc{Cicero.} In that kind, common precepts are given to accuser and defender alike. For whichever of the two shall, by defining and describing the word, have penetrated more deeply into the judge's perception and opinion, and whichever shall have come more nearly and more closely to the common force of the word—the same as the notion of that word which those who hear will have implanted in their minds in rudimentary form—he must win. For this kind is not handled by argument,

\ciceroSection{124} but as it were by unfolding and sifting the word: as, if the accuser should define collusion to be every corruption of a court by the accused—in the case of a man acquitted for a bribe and then summoned back to trial—while the defender defines it not as every corruption, but only as corruption of the accuser by the accused, then let this be first a contest of words; in which, even though the defender's definition should come closer to the usage and intent of speech, the accuser will nonetheless rely on the intent of the law; for he denies that one ought to approve those who framed the laws,

\ciceroSection{125} that, lest a judgement should be held valid if the whole of it had been corrupted, they should not annul it if a single accuser had been corrupted; he takes his stand on equity and advantage, as though a law were now to be written, and he says that they embraced under the single word "collusion" whatever cases of corruption in the courts they then meant to comprehend;

\ciceroSection{126} the defender, on the other hand, calls the usage of speech to witness, and recovers the force of the word from its contrary, as though the accuser were turned about, the contrary name to whom is "collusive defender"; and from the consequences, in that the verdict-letter for the accuser is customarily given to the judge; and from the very name, which signifies the man who in opposite causes seems to stand, as it were, straddled. Yet even he must take refuge in the topics of equity, in the authority of things adjudged, in some limit to be set on the peril. And let this be a common precept: that when each shall have defined, as nearly as he could, to the common sense and force of the word, then he should confirm his own definition and meaning by similar cases and by examples of those who have so spoken.

\ciceroSection{127} And let the accuser, in this kind of case, have that common topic: that it must by no means be conceded that a man who confesses the deed should defend himself by the interpretation of a word; the defender, on the other hand, should rely both on that equity I have set forth, and—when this is on his side—should complain that he is being pressed not by the fact but by a distortion of the word. In this kind he will be able to range over most of the topics of discovery; for he will use both similars and contraries and consequences—although each side will, the defendant, nonetheless, unless his case is plainly absurd, will use them more frequently.

\ciceroSection{128} For the sake of amplification, the things they will wish to say either when they digress from the case or when they sum up—aimed at hostility, or at pity, or at stirring the judges' minds in every way—will be drawn from what has been set forth before, provided only that the magnitude of the matter or of the persons, or odium, or dignity, shall demand it.

\ciceroSection{129} \textsc{Marcus.} I have those points; now I should like to hear what, when it is disputed of what kind a thing is, it is fitting to inquire on either side. \textsc{Cicero.} In that kind, those who are accused confess that they did the very thing for which they are reproached; but since they say they did it lawfully, the whole theory of law must be unfolded by us. And this is divided into two primary parts, nature and statute, and the force of each kind is distributed into divine and human law; of which the one belongs to equity, the other to religion.

\ciceroSection{130} The force of equity, moreover, is twofold: the one part of which is defended on the straightforward principle of the true and the just and, as it is called, the fair and good; the other pertains to the returning of a kindness in turn, which in a benefit is called gratitude, in a wrong, punishment. And these things are common to nature and statute; but proper to statute are both the things which have been written down, and those which are upheld without writing, either by the law of nations or by the custom of our forefathers. Of written law, however, one part is private, another public: public is statute, decree of the Senate, treaty; private is contract, agreement, covenant, stipulation. The things that are not written are upheld either by usage or by the conventions of men and, as it were, by consensus. And this above all, that we should guard our own customs and laws, is in a certain way prescribed by natural law.

\ciceroSection{131} And since the springs, so to speak, of equity have been briefly opened up, we shall have to have prepared, for this kind of case, the things that will have to be said in our speeches concerning nature, concerning laws, concerning the custom of our forefathers, concerning the warding off of wrong, concerning its avenging, concerning every part of law. If a man shall have done something unwittingly, or by necessity, or by chance, which would not be allowed to those who had done it of their own accord and will, then for the excusing of that deed pardon must be sought by craving indulgence, which will be drawn from most of the topics of equity. It has been set forth, as briefly as I could, concerning every kind of dispute, unless you require something further besides.

\ciceroSection{132} \textsc{Marcus.} There is indeed that one thing which I now see remains—of what kind it is, when the dispute turns upon written documents. \textsc{Cicero.} You understand rightly; for once that is set forth, I shall have completed the whole task of my promise. In ambiguity, then, the precepts are common to the two adversaries. For each will defend the sense on which he himself relies as worthy of the writer's prudence; each will defend that the sense which his adversary says is to be understood from the ambiguous writing is either absurd, or useless, or unfair, or base, or that it is even at variance with other writings, whether of others or—above all, if he can manage it—of the same writer; and that the matter and meaning which he himself defends is what any prudent and just man, if a free hand were given him, would have written, but more plainly;

\ciceroSection{133} and that the meaning which he says can be conveyed has in it nothing either of a snare or of fault; whereas, if they should approve the contrary one, the result would be that many faulty, foolish, unfair, and contradictory consequences would follow. But when the writer appears to have meant one thing and written another, the man who relies on the writing will have to set forth the matter and then make use of a reading aloud; thereupon to press his adversary, to repeat, to renew, to ask whether he either denies the writing or disowns having acted against it.

\ciceroSection{134} Afterward let him summon the judge to the force of the written word. Having used this confirmation, let him amplify the matter by praising the law, and let him confound the audacity of one who, although he has openly acted against it and admits it, is nonetheless present and defends the deed. Then let him weaken the defence, when the adversary says that the writer willed one thing, meant another, wrote another: that it must not be tolerated that the framer's intent should be unfolded by anyone rather than by the law itself. Why did he so write, if he did not so mean? Why, when he has neglected the things that are plainly written, does he bring forward things that are nowhere written? Why does he hold that men of the highest prudence in writing are to be condemned for the utmost folly? What hindered the writer from making an exception of that very thing which the adversary says he followed, as though it had been excepted?

\ciceroSection{135} He will use those examples in which the same writer, or—if he cannot find these—in which others, have made the exception of what they thought ought to be excepted. A reason must also be sought, if any can be found, why the exception was not made: it will be said that the law would have been unfair or useless, or that there is one ground for obeying it and another for abrogating it; that the adversary's voice is at odds with the law's. Then, for the sake of amplification, one must speak, with the greatest gravity and vehemence—among other places, but most of all in the peroration—concerning the preservation of the laws, concerning the danger to affairs both public and private.

\ciceroSection{136} The man, on the other hand, who will defend himself by the intent and will of the law, will defend that the force of the law is placed in the design and the mind of the writer, not in the words and the letters; and he will praise the fact that the writer made no exception in the law, lest by-ways should be given to wrongdoing, and so that the judge might interpret the mind of the law from each man's deed. Then he must use examples in which all equity would be thrown into confusion if obedience were given to the words of the laws and not to their meanings.

\ciceroSection{137} Then let this kind of cunning and chicanery be drawn into the judge's hatred, with a certain invidious complaint. And if a case of ignorance shall arise, one that pertains not to wrongdoing but to chance or necessity—a kind we touched on a little before—indulgence will have to be craved by those same equitable considerations against the harshness of the words. But if the writings shall be at variance among themselves, the chain of the art is so great, and most things are so connected and fitted together, that the precepts we gave a little before concerning ambiguity, and those just now concerning the meaning and the writing, may be transferred unchanged to this third kind of case.

\ciceroSection{138} For by the same topics by which, in ambiguity, we defend the sense that aids us, by these same our law must be defended in cases of conflicting laws. Then it must be brought about that we defend the meaning of the one writing by the words of the other. So the very things we just laid down concerning the writing and the meaning, all these we shall transfer here.

\ciceroSection{139} All the subdivisions of oratory have now been set forth for you—the very ones that have flowered from that Academy of ours, and without which they can be neither discovered nor understood nor handled. For to divide a thing itself, and to define, and to distinguish the subdivisions of an ambiguity, and to know the topics of arguments, and to draw the argumentation itself to its conclusion, and to see what is to be assumed in arguing and what follows from the things assumed, and to judge and distinguish the true from the false, the plausible from the incredible, and to censure things ill assumed or ill concluded, and to set out the same matters either narrowly, like those who are called dialecticians, or, as befits an orator, expansively—all this belongs to the art of subtle disputation and copious speech.

\ciceroSection{140} But concerning things good and bad, fair and unfair, useful and useless, honourable and base, what skill or what abundance can the orator have without those arts of the greatest matters? Therefore let these things which I have set forth be to you, my Cicero, as it were the signs of those springs; and if you shall come through to them, whether under these same guides or others, then you will know both these very things better, and many far greater besides. \textsc{Marcus.} I will indeed, and with great zeal, my father; and of all your most splendid gifts I look for none greater.
